\section{结论}

本研究系统地对比了三种深度学习架构在脑肿瘤MRI影像分类任务上的性能，包括纯卷积神经网络（Pure CNN）、纯视觉Transformer（Pure ViT）以及创新性的CNN-Transformer混合模型。通过在包含四类脑肿瘤（胶质瘤、脑膜瘤、无肿瘤、垂体瘤）的公开数据集上进行充分的实验验证，本研究得出以下主要结论：

\textbf{1. 架构互补性得到验证}

实验结果表明，Pure CNN模型凭借其归纳偏置在捕捉局部纹理特征方面表现出色，训练效率高且参数量少；Pure ViT模型通过自注意力机制展现出强大的全局建模能力，能够捕捉远程依赖关系；而混合模型成功融合了两种架构的优势，在整体性能上达到或超过单一架构模型，验证了"局部-全局"特征融合策略的有效性。

\textbf{2. 混合架构的优越性}

本研究提出的CNN-Transformer串行融合架构，通过让CNN模块首先提取局部纹理特征，再由Transformer模块在高层语义特征基础上建立全局依赖关系，实现了计算效率与模型性能的良好平衡。该架构的设计思想与放射科医生"先观察局部细节，再评估全局结构"的诊断流程高度契合，为深度学习在医学影像分析中的应用提供了新的思路。

\textbf{3. 评估体系的全面性}

本研究不仅关注整体准确率，还针对每个类别分别计算了精确率、召回率、F1分数等多种指标，并通过混淆矩阵深入分析了模型的错误模式。这种全面的评估体系更贴近临床实际需求，为模型在真实医疗场景中的应用提供了可靠的性能参考。

\textbf{4. 数据增强与正则化的重要性}

在相对较小的医学影像数据集上，数据增强（随机翻转、旋转、颜色抖动）和正则化技术（Dropout、权重衰减）对于防止过拟合、提升模型泛化能力起到了关键作用。实验过程中采用的余弦退火学习率调度策略也有效促进了模型的稳定收敛。

\textbf{5. 未来研究方向}

尽管本研究取得了良好的实验结果，但仍存在进一步提升的空间。未来研究可以从以下几个方向展开：
\begin{itemize}
    \item 引入大规模医学影像预训练权重，提升模型在小规模数据集上的性能；
    \item 探索注意力可视化技术，增强模型的可解释性，提升临床医生的信任度；
    \item 开发多模态融合模型，整合不同MRI序列（T1、T2、FLAIR）的互补信息；
    \item 引入不确定性量化方法，为预测结果提供置信度估计，辅助临床决策；
    \item 探索模型轻量化技术，使其能够部署到移动端或边缘设备，拓展应用场景。
\end{itemize}

\textbf{6. 临床应用价值}

本研究开发的混合模型具有重要的临床应用价值。它不仅可以作为放射科医生的辅助诊断工具，减轻其工作负担、降低误诊和漏诊风险，还可以用于基层医疗机构的快速初筛，以及医学生和低年资医生的培训。然而，必须强调的是，深度学习模型应始终作为辅助工具而非替代专家判断，最终诊断决策权应由专业医生掌握，以确保患者安全。

总而言之，本研究通过系统的实验设计和深入的分析，验证了CNN与Transformer融合架构在脑肿瘤MRI影像分类任务上的优越性能，为深度学习在医学影像分析领域的应用提供了有价值的理论支撑和实践经验。随着医学影像数据的不断积累和深度学习技术的持续进步，我们有理由相信，人工智能将在辅助医生进行更精准、高效的疾病诊断方面发挥越来越重要的作用。
